\chapter{Giới thiệu}
\label{chap:introduction}

Các ứng dụng web hiện nay là nền tảng của thương mại điện tử, dịch vụ công, giáo dục trực tuyến, ngân hàng số và nhiều hệ thống nghiệp vụ. Khi phần lớn dữ liệu và quy trình vận hành được đưa lên môi trường web, lớp ứng dụng trở thành mục tiêu tấn công quan trọng. Trong các rủi ro phổ biến, Cross-Site Scripting (XSS) và SQL Injection (SQLi) vẫn có ảnh hưởng lớn vì chúng khai thác trực tiếp cách ứng dụng xử lý dữ liệu đầu vào của người dùng~\cite{owasp2021top10}. XSS có thể dẫn tới đánh cắp phiên đăng nhập, thực thi mã JavaScript độc hại trên trình duyệt nạn nhân hoặc thay đổi nội dung hiển thị. SQL Injection có thể cho phép kẻ tấn công truy vấn, sửa đổi hoặc phá hoại cơ sở dữ liệu phía sau ứng dụng.

Một hướng phòng thủ phổ biến là triển khai Web Application Firewall (WAF) trước ứng dụng. WAF không thay thế việc lập trình an toàn, nhưng đóng vai trò như một lớp kiểm tra bổ sung: mọi HTTP request đi qua WAF sẽ được phân tích để phát hiện các mẫu tấn công trước khi tới ứng dụng. Trong hệ sinh thái mã nguồn mở, ModSecurity là một WAF engine được sử dụng rộng rãi, còn OWASP Core Rule Set (CRS) là bộ luật phổ biến được triển khai cùng ModSecurity để phát hiện nhiều nhóm tấn công web khác nhau~\cite{modsecurity-crs-overview,owasp-crs-docs}. Tuy nhiên, triển khai WAF hiệu quả không chỉ là bật càng nhiều rule càng tốt. Một hệ thống WAF thực tế cần cân bằng giữa hai mục tiêu: phát hiện được tấn công thật và hạn chế chặn nhầm request hợp lệ.

\section{Bối cảnh an toàn ứng dụng web}
\label{sec:motivation}

Về bản chất, WAF là một cơ chế phòng thủ dựa trên luật. Mỗi request được so khớp với một tập quy tắc mô tả dấu hiệu tấn công. Ví dụ, một request chứa chuỗi \path|<script>alert(1)</script>| có thể bị xem là XSS, hoặc một tham số chứa \path|' OR 1=1--| có thể bị xem là SQL Injection. Nếu chỉ nhìn ở mức trực giác, cách làm này có vẻ đơn giản: càng liệt kê nhiều dấu hiệu nguy hiểm thì WAF càng phát hiện tốt. Tuy nhiên, trong thực tế, cùng một từ khóa hoặc ký tự có thể xuất hiện trong cả ngữ cảnh tấn công và ngữ cảnh hợp lệ. Một website có thể truyền đoạn JavaScript hợp lệ trong tham số, một nội dung hướng dẫn lập trình có thể chứa từ \path|SELECT|, hoặc một form HTML hợp lệ có thể dùng event handler như \path|onsubmit="validate()"|.

Trong bối cảnh nghiên cứu này, vấn đề trọng tâm không phải là xây dựng một WAF hoàn toàn mới, mà là tối ưu bộ luật phát hiện XSS và SQLi trong môi trường ModSecurity/OWASP CRS. Đề tài tập trung vào câu hỏi: làm thế nào để các rule vẫn phát hiện được các kỹ thuật bypass phổ biến, nhưng không chặn nhầm quá nhiều request hợp lệ? Đây là bài toán quan trọng vì false positive cao có thể khiến WAF khó triển khai trong thực tế. Nếu WAF liên tục chặn nhầm người dùng thật, quản trị viên thường buộc phải tắt rule hoặc tăng ngưỡng chặn, từ đó làm giảm hiệu quả bảo vệ.

\section{Thực trạng WAF rule-based và vấn đề nghiên cứu}
\label{sec:problem-overview}

OWASP CRS là một nền tảng quan trọng vì cung cấp sẵn nhiều luật phát hiện tấn công web phổ biến. CRS cũng dùng cơ chế cộng điểm bất thường: mỗi dấu hiệu đáng ngờ làm request tăng điểm, và request chỉ bị chặn khi tổng điểm vượt ngưỡng cấu hình. Cách làm này giúp WAF kết hợp nhiều tín hiệu khác nhau thay vì phụ thuộc vào một rule đơn lẻ. Tuy nhiên, khi cấu hình nghiêm ngặt hơn hoặc bật nhiều rule hơn, WAF thường phát hiện được nhiều dấu hiệu hơn nhưng cũng dễ chặn nhầm request hợp lệ hơn~\cite{owasp-crs-annotated,paranoia-levels-study,waf-benchmark-2024}.

Một hướng cải thiện phổ biến là bổ sung custom rules cho các nhóm tấn công cần ưu tiên, chẳng hạn XSS và SQL Injection. Các nghiên cứu trước cho thấy việc bổ sung rule có thể cải thiện khả năng phát hiện trên benchmark~\cite{researcher-group-2025}. Tuy nhiên, nhiều rule kiểu này vẫn dựa nhiều vào cách liệt kê dấu hiệu đã biết: liệt kê tên hàm JavaScript, tag HTML, từ khóa SQL hoặc một số payload cụ thể. Cách tiếp cận này dễ hiểu và có hiệu quả với các mẫu đã biết, nhưng có ba vấn đề chính ở mức tổng quan.

Thứ nhất, danh sách liệt kê khó bao phủ hết các biến thể tấn công. Kẻ tấn công có thể thay đổi cú pháp, mã hóa payload, chèn khoảng trắng bất thường hoặc dùng cách gọi hàm khác để né rule. Thứ hai, rule quá rộng có thể chặn nhầm dữ liệu hợp lệ, vì người dùng có thể nhập nội dung kỹ thuật chứa HTML, JavaScript hoặc SQL mà không có ý định tấn công. Thứ ba, khi danh sách payload ngày càng dài, rule trở nên khó đọc, khó kiểm thử và khó bảo trì.

Từ thực trạng trên, khóa luận lựa chọn hướng \textbf{generalized rules}: thay vì chỉ hỏi ``request có chứa từ khóa nguy hiểm hay không'', rule cần xét xem dữ liệu đó có nằm trong ngữ cảnh nguy hiểm hay không. Với XSS, trọng tâm là ngữ cảnh có thể thực thi JavaScript. Với SQLi, trọng tâm là cấu trúc có khả năng làm thay đổi ý nghĩa truy vấn. Chương 2 sẽ trình bày chi tiết hơn cơ sở lý thuyết, các phương pháp liên quan và hạn chế của từng hướng; Chương 3 trình bày giải pháp đề xuất dựa trên các phân tích đó.

\section{Mục tiêu nghiên cứu}
\label{sec:objectives}

Mục tiêu tổng quát của khóa luận là \textbf{đánh giá và đề xuất một hướng tối ưu bộ luật WAF ModSecurity cho phát hiện XSS và SQL Injection nhằm giảm false positives, duy trì khả năng phát hiện cao trên các payload tấn công phổ biến, và cải thiện tính bảo trì so với phương pháp enumeration-based}. Mục tiêu này nhấn mạnh kết quả cần đạt được của nghiên cứu, không chỉ dừng ở việc triển khai thêm một tập rule mới.

Các mục tiêu cụ thể bao gồm:

\begin{enumerate}
    \item \textbf{Làm rõ hạn chế của CRS và rule enumeration-based} trong phát hiện XSS/SQLi, đặc biệt về false positives, coverage trước payload biến thể và khả năng bảo trì.

    \item \textbf{Đánh giá khả năng giảm false positives của hướng generalized/context-aware detection}, trong đó rule cần xét ngữ cảnh thực thi hoặc cấu trúc tấn công thay vì bắt từ khóa đơn lẻ.

    \item \textbf{Đánh giá vai trò của decode pipeline} đối với payload encoded/obfuscated, bao gồm URL encoding, HTML entity, JavaScript escape, null byte, whitespace obfuscation và Base64.

    \item \textbf{Đánh giá khả năng mở rộng coverage SQLi} khi sử dụng semantic patterns trên nhiều nhóm kỹ thuật và nhiều hệ quản trị cơ sở dữ liệu, thay vì chỉ bắt một số payload cố định.

    \item \textbf{Đánh giá tác động của tuning exclusions} tới false positive rate và recall trong phạm vi XSS/SQLi, đồng thời ghi nhận các trade-off bảo mật khi vô hiệu hóa một số CRS rules.

    \item \textbf{Phân tích ưu và nhược điểm của WAF rule-based sau tuning} khi đặt cạnh hướng AI-based detection về khả năng giải thích, độ trễ, khả năng cập nhật và tính phù hợp với triển khai real-time.
\end{enumerate}

\section{Câu hỏi nghiên cứu}
\label{sec:research-questions}

Từ mục tiêu trên, khóa luận tập trung trả lời các câu hỏi nghiên cứu sau:

\begin{itemize}
    \item \textbf{RQ1:} Phương pháp generalized/context-aware rules có giúp giảm false positives so với phương pháp enumeration-based trong phát hiện XSS hay không?

    \item \textbf{RQ2:} Việc sử dụng multi-decode pipeline có cải thiện khả năng phát hiện các payload XSS/SQLi được mã hóa hoặc obfuscate hay không?

    \item \textbf{RQ3:} Các SQLi rules dựa trên semantic patterns có mở rộng coverage tốt hơn so với rule chỉ bắt một số payload cố định hay không?

    \item \textbf{RQ4:} Việc tuning CRS bằng cách vô hiệu hóa có chọn lọc các rule gây FP cao có thể giảm false positives mà vẫn giữ recall cao cho XSS/SQLi hay không?

    \item \textbf{RQ5:} Khi so sánh với hướng AI-based detection, WAF rule-based sau tuning có những ưu và nhược điểm gì về khả năng giải thích, độ trễ, khả năng cập nhật và tính phù hợp với triển khai real-time?
\end{itemize}

Các câu hỏi này định hướng cấu trúc của các chương sau. Chương 2 phân tích cơ sở lý thuyết, nghiên cứu liên quan và hạn chế của các phương pháp hiện có; Chương 3 trình bày phương pháp nghiên cứu và giải pháp đề xuất; Chương 4 trình bày thiết kế thực nghiệm và phương pháp đánh giá; Chương 5 trình bày kết quả thực nghiệm; và Chương 6 thảo luận ý nghĩa, giới hạn và trade-off của các kết quả đó.

\section{Phạm vi nghiên cứu}
\label{sec:scope}

Để đảm bảo bài toán đủ tập trung, khóa luận giới hạn phạm vi như sau:

\begin{itemize}
    \item \textbf{Loại tấn công chính:} XSS và SQL Injection. Một số pattern NoSQL hoặc OS-command-via-SQL có thể xuất hiện trong lớp heuristic, nhưng không phải trọng tâm chính.

    \item \textbf{Nền tảng triển khai:} ModSecurity kết hợp OWASP CRS và custom rules. Đề tài không xây dựng WAF engine mới.

    \item \textbf{Dữ liệu đánh giá:} các dataset và benchmark payloads như GoTestWAF, Mereani XSS, XSSed-DMOZ, CSIC 2010 và ECML nếu đã được chuẩn hóa trong pipeline thực nghiệm.

    \item \textbf{Phương pháp đánh giá:} tập trung vào TP, TN, FP, FN và các metrics như Recall, Precision, FPR, Accuracy và F1-score.

    \item \textbf{Tuning:} tập trung vào giảm false positives trong phạm vi XSS/SQLi bằng cách phân tích benchmark và audit logs.
\end{itemize}

Các nội dung ngoài phạm vi gồm: xem RCE/Shell Injection là mục tiêu chính; triển khai production đầy đủ; thay thế toàn bộ OWASP CRS; và cam kết phát hiện mọi zero-day hoặc mọi biến thể tấn công mới. Những giới hạn và rủi ro liên quan tới tuning sẽ được thảo luận chi tiết hơn ở Chương 6.

\section{Đóng góp chính của khóa luận}
\label{sec:contributions}

Khóa luận có các đóng góp chính sau:

\begin{enumerate}
    \item \textbf{Kiến trúc phát hiện XSS context-aware theo nhiều lớp.} Bộ XSS rules v2 sử dụng dải ID \path|9942xxx|, tập trung vào execution context thay vì chỉ bắt từ khóa. Các lớp gồm core XSS, structural injection, encoded/obfuscated payload, DOM context và advanced heuristic.

    \item \textbf{Bộ SQLi rules v2 dựa trên semantic patterns.} Bộ SQLi rules sử dụng dải ID \path|9943xxx|, phát hiện các nhóm kỹ thuật như UNION-based, stacked queries, boolean/time-based blind SQLi, error-based SQLi, OOB SQLi, schema enumeration, comment obfuscation và encoded SQL keywords.

    \item \textbf{Multi-decode pipeline cho payload bị mã hóa.} Giải pháp áp dụng các transform như \path|t:urlDecodeUni|, \path|t:htmlEntityDecode|, \path|t:jsDecode|, \path|t:removeNulls| và \path|t:compressWhitespace|. Ngoài ra, một lớp riêng xử lý Base64 được sử dụng để phát hiện payload sau khi decode.

    \item \textbf{Phương pháp tuning dựa trên dữ liệu.} Khóa luận sử dụng benchmark và audit logs để xác định rule gây FP cao, sau đó phân nhóm rủi ro và cấu hình \path|TUNING-EXCLUSIONS.conf| có kiểm soát.

    \item \textbf{Benchmark framework và đánh giá thực nghiệm.} Hệ thống benchmark tự động gửi payload tới WAF, phân loại TP/TN/FP/FN và sinh báo cáo để so sánh rule gốc, CRS baseline, custom generalized rules v2 và AI model tham chiếu.

    \item \textbf{Phân tích trade-off minh bạch.} Khóa luận không chỉ trình bày kết quả tốt, mà còn thảo luận rủi ro của tuning, đặc biệt là việc vô hiệu hóa một số CRS rules ngoài phạm vi XSS/SQLi như RCE/Shell và trade-off khi tắt libinjection SQLi.
\end{enumerate}

\section{Cấu trúc khóa luận}
\label{sec:thesis-structure}

Phần còn lại của khóa luận gồm năm chương. Chương 2 trình bày cơ sở lý thuyết và các nghiên cứu liên quan. Chương 3 trình bày phương pháp nghiên cứu và giải pháp đề xuất. Chương 4 mô tả thiết kế thực nghiệm và phương pháp đánh giá. Chương 5 trình bày kết quả thực nghiệm và đánh giá. Chương 6 thảo luận các kết quả chính, nêu giới hạn và kết luận khóa luận.
